%----------------------------------------
\begin{frame}
  \frametitle{Part 5}

{\Large\bf How \spin{} works}

\end{frame}

%----------------------------------------
\begin{frame}
  \frametitle{The automata-theoretic approach} \small

\begin{enumerate}
\item The \promelaname{} model is an automaton $A_M$ whose language
      $\mathcal{L}(A_M)$ is the set of all its executions.
\pause
\item The negated property $\neg\varphi$ is translated into a \buchi{} automaton
      $A_{\neg\varphi}$ accepting exactly the {\bf violating} executions.
\pause
\item \spin{} builds the {\bf synchronous product} $A_M \otimes A_{\neg\varphi}$.
\pause
\item Then:
\[
M \models \varphi
\quad\Longleftrightarrow\quad
\mathcal{L}(A_M \otimes A_{\neg\varphi}) = \emptyset
\]
\end{enumerate}

\pause
\vspace{0.4cm}

\begin{block}{}\small
Verification $=$ {\bf emptiness checking}. And a \buchi{} automaton is non-empty
iff some accepting state is reachable {\em and} lies on a cycle.\\[0.2cm]
So: {\bf safety} $=$ find a bad state. {\bf Liveness} $=$ find a bad cycle.
\end{block}

\end{frame}

%----------------------------------------
\begin{frame}
  \frametitle{The search} \small

\begin{block}{Safety: depth-first reachability}\small
Push the initial state; expand successors; hash each into a table; report an error
the first time an assertion fails, an end state is invalid, or a deadlock occurs.
The DFS stack {\bf is} the counterexample.
\end{block}

\pause

\begin{block}{Liveness: a second search, hunting for a loop}\small
A liveness violation is an {\em infinite} bad run --- and in a finite model, an
infinite run means a {\bf loop}. The property automaton marks a state
{\em accepting} when the violation is under way and could go on forever --- for
\ipromela@[]<> g0@ that is ``\ipromela{g0} has not held since''.
So search as usual, and each time you finish exploring an accepting state, start a
{\bf second} search from it asking one question: {\em can I get back here?} If it
can, that loop {\em is} the counterexample.\\[0.15cm]
The {\em nested} depth-first search: linear time, two bits per state --- seen by
the first, seen by the second.
\end{block}

\pause
\vspace{0.15cm}

{\bf On-the-fly:} \spin{} never builds the state graph first --- it constructs
states as it searches and stops at the first violation. A bug in a
$10^{20}$-state system can surface after $10^3$ states.

\end{frame}

%----------------------------------------
\begin{frame}
  \frametitle{The enemy: state explosion} \small

The number of states is roughly

\[
|S| \;\approx\; \prod_{i} |P_i| \;\times\; \prod_{j} |V_j| \;\times\; \prod_{k} |C_k|
\]

--- process locations $\times$ variable valuations $\times$ channel contents.

\pause
\vspace{0.3cm}

{\bf Everything is multiplicative.} One extra \ipromela{byte} variable multiplies
the state space by up to 256. One extra process multiplies it by that process's
entire local state space.

\pause
\vspace{0.4cm}

\begin{block}{Practical implication}\small
Model-checking effort is dominated not by writing \promelaname{} but by
{\bf choosing what to leave out}.
\end{block}

\end{frame}

%----------------------------------------
\begin{frame}
  \frametitle{\spin{}'s weapons against state explosion} \small

\begin{itemize}
\item {\bf Skip the orders that cannot matter} ({\em partial-order reduction}, on
      by default). If two processes are doing things that cannot affect each
      other, try one order and not the other. Usually an order of magnitude, free.
\item {\bf Pack each state tighter} (\tool{-DCOLLAPSE}). Slower, smaller, exact.
\item {\bf Do not keep the states --- keep hashes of them.} One idea at two
      settings; either can miss a bug, neither can invent one.
  \begin{itemize}\small
  \item a {\bf fingerprint} per state (\tool{-DHC}): a few bytes each, so memory
        still grows with the number of states. Errs only if two states collide.
  \item {\bf no entries at all} (\tool{-DBITSTATE}): one fixed-size scoreboard ---
        tick a few boxes per state, ``already ticked'' means ``been here''. A few
        bits rather than 384, whatever the state's size. Errs once other states
        have ticked those boxes between them.
  \end{itemize}
\item {\bf Search in parallel} (\tool{-DNCORE=N}, {\em swarm}): many short searches,
      each looking elsewhere. Good at finding bugs; not a proof.
\item {\bf Give up past a chosen depth} (\tool{-mN}). If the run reports
      \tool{depth reached} equal to your bound it never finished, and
      \tool{errors: 0} means nothing.
\end{itemize}

\pause

{\bf Rule of thumb:} exhaustive if it fits; otherwise a lossy search that finds the
bug beats an exhaustive one that runs out of memory.

\end{frame}
